% =====================================================================
%  CHƯƠNG 3. DỮ LIỆU                              | Phụ trách: Cả 2
%  Ngân sách: mục tiêu 9 trang, trần 10 trang
%  Nguồn: docs/research/quy-dinh-bien-so-xe-vn.md            (mục 3.1),
%         src/ml/notebooks/eda-plate-datasets.ipynb           (mục 3.2),
%         docs/research/eda_outputs/eda_summary.csv           (Bảng 3.4),
%         data/processed/a1_det/phash_report.txt              (mục 3.5)
% =====================================================================
\chapter{Dữ Liệu}
\label{ch:dulieu}

% =====================================================================
\section{Quy Định Biển Số Xe Việt Nam}
\label{sec:quydinhbienso}
% =====================================================================

Cấu trúc pháp lý của biển số quyết định cách thiết kế nhãn dữ liệu, bộ ký tự OCR và quy
tắc hậu xử lý.

\subsection{Căn Cứ Pháp Lý Hiện Hành}

Quy định biển số hiện hành (từ 01/01/2025) gồm Luật Trật tự, an toàn giao thông đường bộ
2024, lần đầu đưa màu nền biển vào luật, cùng ba Thông tư của Bộ Công an: 79/2024 về bộ
seri, số ký tự và kích thước theo loại xe (thay Thông tư 58/2020), 81/2024 ban hành quy
chuẩn QCVN 08:2024/BCA về kích thước và chất liệu biển, và 24/2023 quản lý biển theo mã
định danh chủ xe (biển gắn theo chủ, không theo
xe)~\autocite{luatttatgt2024,tt792024bca,tt812024bca,tt242023bca}.

\subsection{Cấu Trúc Chuỗi Ký Tự}

Định dạng tổng quát của biển số xe cơ giới Việt Nam là:

\begin{center}
\texttt{[Mã tỉnh 2 chữ số][gạch ngang][Seri 1 chữ cái] [Số thứ tự 5 chữ số, nhóm 3.2]}
\end{center}

\noindent Ví dụ biển một dòng (ô tô) \texttt{30A-123.45}; biển hai dòng (xe máy) dòng
trên \texttt{59N}, dòng dưới \texttt{187.51}. Ba thành phần có hệ quả trực tiếp với thiết
kế OCR:

\begin{itemize}
  \item \textbf{Mã tỉnh}: 2 chữ số từ 11 đến 99. Từ 01/07/2025, cả nước sáp nhập từ 63
        xuống 34 tỉnh thành (Thông tư 51/2025/TT-BCA), nhưng xe đăng ký trước đó giữ biển
        cũ, nên nhiều mã tỉnh cũ vẫn lưu hành song song nhiều năm.
  \item \textbf{Seri}: chọn trong tập \textbf{20 chữ cái} (A, B, C, D, E, F, G, H, K, L,
        M, N, P, S, T, U, V, X, Y, Z); sáu chữ I, J, O, Q, R, W không dùng cho biển dân
        sự. Bộ 20 chữ cái này được dùng để ràng buộc hậu xử lý OCR, loại nhầm lẫn O với 0
        và I với 1.
  \item \textbf{Số thứ tự}: 5 chữ số nhóm 3.2, từ 000.01 đến 999.99. Riêng seri MĐ và TĐ
        của xe máy điện chứa thêm ký tự ``Đ''.
\end{itemize}

\subsection{Bố Cục Biển Theo Loại Xe}

Biển dài một dòng và biển ngắn hai dòng không phải hai kiểu thay thế nhau: với ô tô, đây
là hai biển bắt buộc gắn đồng thời trên cùng một xe (Bảng~\ref{tab:bienso-bocuc}).

\begin{table}[htbp]
\centering
\caption{Bố cục biển số theo loại phương tiện.}

{\footnotesize Nguồn: \autocite{tt792024bca,tt812024bca}.}
\label{tab:bienso-bocuc}
\small
\begin{tabular}{p{3.2cm} p{1.6cm} p{3.0cm} p{3.0cm}}
\toprule
\textbf{Loại xe} & \textbf{Số biển} & \textbf{Bố cục} & \textbf{Kích thước (mm)} \\
\midrule
Ô tô (con, tải, khách) & 2 biển & Biển trước: 1 dòng 520$\times$110; biển sau: 2 dòng 330$\times$165 & 2 kích thước khác nhau \\
\addlinespace
Mô tô, xe gắn máy (kể cả xe máy điện) & 1 biển (sau) & Luôn 2 dòng & 190$\times$140 \\
\addlinespace
Máy kéo, rơ-moóc & 1 biển (sau) & 2 dòng & 330$\times$165 \\
\bottomrule
\end{tabular}
\end{table}

Hệ quả cho pipeline: camera cổng có thể thu biển một dòng hoặc hai dòng của cùng một ô tô
tùy góc đặt (trước hay sau xe), nên dữ liệu huấn luyện cần cân bằng cả hai bố cục.

\subsection{Màu Nền Biển và Ý Nghĩa Phân Loại Phí}
\label{sec:maunenbien}

Màu nền biển phản ánh nhóm chủ sở hữu theo pháp luật và là cơ sở áp biểu phí
(Chương~\ref{ch:hethong}). Bảng~\ref{tab:maubienso} tổng hợp bốn màu chính theo Luật Trật
tự, an toàn giao thông đường bộ 2024~\autocite{luatttatgt2024}.

\begin{table}[htbp]
\centering
\caption{Màu nền biển số xe và nhóm chủ sở hữu theo quy định hiện hành.}

{\footnotesize Nguồn: \autocite{luatttatgt2024,tt792024bca}.}
\label{tab:maubienso}
\begin{tabular}{p{2.2cm} p{4.8cm} p{3.6cm}}
\toprule
\textbf{Màu nền} & \textbf{Đối tượng} & \textbf{Ưu tiên đề tài} \\
\midrule
Trắng (chữ đen) & Cá nhân, tổ chức dân sự trong nước & Bắt buộc (chiếm đa số) \\
Vàng (chữ đen) & Xe kinh doanh vận tải: taxi, xe khách, xe tải kinh doanh & Bắt buộc \\
Xanh dương (chữ trắng) & Cơ quan Đảng, Nhà nước, đơn vị sự nghiệp công lập & Tùy chọn \\
Đỏ (chữ trắng) & Xe quân sự (Bộ Quốc phòng) & Tùy chọn \\
\bottomrule
\end{tabular}
\end{table}

Xe điện không có màu biển riêng, chỉ dán thêm tem nhận diện màu xanh lá. Màu thực tế trên
ảnh camera có thể lệch so với chuẩn CIE, đặc biệt dưới đèn flash hồng ngoại ban đêm, nên
cần xử lý CLAHE trước khi phân loại màu biển (Chương~\ref{ch:ocr}).

% =====================================================================
\section{Khảo Sát Bộ Dữ Liệu Công Khai}
\label{sec:khaosatdataset}
% =====================================================================

Đề tài khảo sát các bộ dữ liệu công khai theo hai nhóm: (A) biển số Việt Nam
(Bảng~\ref{tab:dataset-a}) và (B) phương tiện và bối cảnh Việt Nam
(Bảng~\ref{tab:dataset-b}); mỗi bộ được kiểm tra trực tiếp trên trang nguồn trong tháng 7
năm 2026. Các bộ ALPR và bãi xe nước ngoài như CCPD, PKLot, CNRPark+EXT, AOLP và OpenALPR
benchmarks~\autocite{xu2018ccpd,almeida2015pklot,amato2017cnrpark,hsu2013aolp,openalpr2016benchmarks}
chỉ dùng để đối chiếu vì không có biển số Việt Nam.

\subsection{Nhóm A: Biển Số Việt Nam}

\begin{table}[htbp]
\centering
\caption{Các bộ dữ liệu nhóm A: biển số xe Việt Nam.}

{\footnotesize Nguồn: Kaggle, Roboflow Universe.}
\label{tab:dataset-a}
\small
\begin{tabular}{p{0.8cm} p{3.2cm} p{1.4cm} p{1.6cm} p{1.6cm} p{2.0cm}}
\toprule
\textbf{Ký hiệu} & \textbf{Tên bộ dữ liệu} & \textbf{Host} & \textbf{Số ảnh} & \textbf{Giấy phép} & \textbf{Biển 1 hàng/2 hàng} \\
\midrule
A1 & Vietnam License Plate Segment (duydieunguyen)~\autocite{duydieunguyen2023vnplatesegment} & Kaggle & khoảng 5.000 & Không rõ & Có, tách rõ số lượng \\
A2 & Vietnamese Car License Plate (Cuong Ta)~\autocite{cuongta2021vnplate} & Roboflow & 8.255 & CC0 1.0 & Không rõ tỷ lệ \\
A3 & vietnamese-license-plate-tptd0 (school-fuhih)~\autocite{schoolfuhih2023vnplate} & Roboflow & 8.397 & CC BY 4.0 & Không rõ tỷ lệ \\
A4 & Motorcycle License Plate (HaUI)~\autocite{haui2023motorcycleplate} & Roboflow & 1.748 & MIT & Thiên về xe máy (2 hàng) \\
A5 & greenParking (thaiphung) & Roboflow & 1.748 & CC BY 4.0 & Không xác định \\
A6 & Viet Nam OCR plate (license-plate-reg)~\autocite{licenseplatereg2025ocrplate} & Roboflow & 3.819 & CC0 1.0 & Không rõ tỷ lệ \\
A7 & VNLicensePlate\_yolov7 (bomaich)~\autocite{bomaich2022vnlicenseplate7} & Kaggle & 1.000 & Không rõ & Có, không tách số \\
\bottomrule
\end{tabular}
\end{table}

Không bộ nào trong nhóm A ghi rõ điều kiện ánh sáng hay bối cảnh chụp.

\subsection{Nhóm B: Phương Tiện và Bối Cảnh Đường Phố Việt Nam}

\begin{table}[htbp]
\centering
\caption{Các bộ dữ liệu nhóm B: phương tiện và bối cảnh Việt Nam.}

{\footnotesize Nguồn: Roboflow Universe, GitHub.}
\label{tab:dataset-b}
\small
\begin{tabular}{p{0.8cm} p{3.4cm} p{1.4cm} p{1.3cm} p{1.4cm} p{2.5cm}}
\toprule
\textbf{Ký hiệu} & \textbf{Tên bộ dữ liệu} & \textbf{Host} & \textbf{Số ảnh} & \textbf{Giấy phép} & \textbf{Ghi chú} \\
\midrule
B1 & Vietnamese vehicle v3 (car-classification)~\autocite{carclassification2023vnvehicle} & Roboflow & 1.547 & CC BY 4.0 & 4 lớp gốc car/bus/truck/motorcycle, đã được ánh xạ lại thành 8 lớp \\
B2 & Vehicle Vietnam-CanTho~\autocite{vehicle2022cantho} & Roboflow & 1.110 đến 1.746 & CC BY 4.0 & Chỉ ban ngày, đường phố Cần Thơ \\
B3 & UIT-CVID21~\autocite{uitcvid21} & GitHub (UIT) & 10.000 (chưa xem được thật) & Không rõ & Chụp từ drone, khác góc camera cổng bãi xe \\
B4 & QuangTranUTE/Vehicle-Detection & GitHub & khoảng 11.000 & Không rõ & Camera giám sát giao thông, không phải bãi xe \\
B5 & Vehicle Body Style Dataset~\autocite{researchprojects2022vehiclebodystyle} & Roboflow & 10.000 & CC BY 4.0 & 12 kiểu dáng theo phân loại CompCars, gốc Trung Quốc \\
\bottomrule
\end{tabular}
\end{table}

Chưa bộ dữ liệu Việt Nam nào ở nhóm B tách phương tiện theo kiểu dáng thân xe (sedan, SUV,
pickup), nên đề tài bổ sung B5. B5 có 12 lớp gần cân bằng (lớp lớn nhất chỉ gấp 1,02 lần lớp nhỏ
nhất), nhưng xe chiếm trung bình 0,31 diện tích ảnh và không có xe máy, phản ánh phong cách
chụp tại đại lý, khác ảnh camera cổng; hạn chế này được phân tích ở
Chương~\ref{ch:phanloai}.

\subsection{Đánh Giá và Hướng Lựa Chọn}
\label{sec:chondataset}

Đề tài chọn A1 làm bộ chính cho phát hiện biển vì là bộ duy nhất tách rõ biển một hàng
(BSD) và biển hai hàng (BSV). A1 cùng A3 và A4 được tải về và phân tích khám phá ở
Mục~\ref{sec:edabienso} để xác minh bối cảnh chụp.

% =====================================================================
\section{Phân Tích Khám Phá Dữ Liệu Biển Số}
\label{sec:edabienso}
% =====================================================================

Ba bộ A1, A3, A4 được tải và phân tích khám phá độc lập, kết quả tổng hợp ở
Bảng~\ref{tab:eda-summary}.

\subsection{Tổng Quan Số Liệu}

\begin{table}[htbp]
\centering
\caption{Tổng hợp kết quả EDA trên ba bộ dữ liệu biển số được chọn.}

{\footnotesize Độ sáng đo trên mẫu ngẫu nhiên $n = 600$ ảnh mỗi bộ, seed cố định 42.}
\label{tab:eda-summary}
\begin{tabular}{lrrp{3.4cm}rr}
\toprule
\textbf{Bộ dữ liệu} & \textbf{Số ảnh} & \textbf{Số object} & \textbf{Độ phân giải (px)} & \textbf{Độ sáng TB} & \textbf{\% ảnh tối} \\
\midrule
A3 & 8.357 & 8.548 & 640$\times$640 (cố định) & 100,7 & 12,8\% \\
A4 & 1.748 & 1.748 & 472$\times$303 (cố định) & 111,1 & 0,0\% \\
A1 & 4.578 & 5.200 & 380 đến 4032 $\times$ 285 đến 3024 & 102,1 & 10,7\% \\
\bottomrule
\end{tabular}
\end{table}

\subsection{Phân Bố Lớp và Kích Thước Biển Số}

Hình~\ref{fig:eda-class} cho thấy ở A1, lớp BSV (biển hai hàng, 3.559 đối tượng) nhiều gấp
khoảng 2,2 lần lớp BSD (biển một hàng, 1.641 đối tượng); nếu gộp ba bộ để huấn luyện cần
xử lý mất cân bằng này để mô hình không thiên về biển hai hàng.

\begin{figure}[htbp]
\centering
\includegraphics[width=0.90\textwidth]{eda_class_distribution.png}
\caption{Phân bố số lượng object gán nhãn theo bộ dữ liệu và theo lớp (BSD và BSV).}

\label{fig:eda-class}
\end{figure}

Tỷ lệ diện tích vùng biển trên toàn ảnh trung bình chỉ khoảng
0,028 đến 0,043 ở cả ba bộ (Hình~\ref{fig:eda-arearatio}): phần lớn ảnh chụp toàn cảnh đầu
xe chứ không cận cảnh sát biển, đúng đặc thù camera cổng bãi xe.

\begin{figure}[htbp]
\centering
\includegraphics[width=0.90\textwidth]{eda_plate_area_ratio.png}
\caption{Phân bố tỷ lệ diện tích biển số trên diện tích ảnh cho ba bộ dữ liệu.}

\label{fig:eda-arearatio}
\end{figure}

\subsection{Kích Thước Ảnh và Độ Sáng}

A3 và A4 đã bị Roboflow resize về kích thước cố định khi xuất (640$\times$640 và
472$\times$303, độ lệch chuẩn bằng không) nên độ phân giải gốc coi như mất; chỉ A1 tải
trực tiếp từ Kaggle còn giữ kích thước gốc, rộng từ 380 đến 4032 pixel
(Hình~\ref{fig:eda-size}). Do không bộ nào có nhãn ánh sáng, độ sáng trung bình ảnh xám
(thang 0 đến 255) được dùng làm proxy tạm: A4 gần như không biến thiên, A3 và A1 có khoảng
11 đến 13\% ảnh ở mức tối.

\begin{figure}[htbp]
\centering
\includegraphics[width=0.90\textwidth]{eda_size_brightness.png}
\caption{Phân bố kích thước ảnh, tỷ lệ khung hình và độ sáng trên ba bộ dữ liệu.}

\label{fig:eda-size}
\end{figure}

\subsection{Xem Mẫu Ảnh để Xác Minh Bối Cảnh}

Ảnh mẫu (Hình~\ref{fig:eda-sample}) cho thấy cả ba bộ đều chụp từ camera cổng bãi xe hoặc
hầm gửi xe, có thanh chắn trong khung, nhiều ảnh có timestamp kiểu camera DVR và ảnh hầm
thiếu sáng, đúng bối cảnh đề tài.

\begin{figure}[htbp]
\centering
\includegraphics[width=0.90\textwidth]{eda_sample_context.png}
\caption{Ảnh mẫu ngẫu nhiên từ ba bộ dữ liệu, minh họa bối cảnh camera cổng bãi xe.}

\label{fig:eda-sample}
\end{figure}

% =====================================================================
\section{Bộ Dữ Liệu Bổ Sung cho OCR, Màu Biển và Phân Loại Xe}
\label{sec:datasetocrcolor}
% =====================================================================

\subsection{Bộ Dữ Liệu Phát Hiện Biển và OCR}

Bảng~\ref{tab:dataset-ocr} liệt kê các bộ đã tải thực tế phục vụ phát hiện biển, OCR và ký
tự đơn lẻ.

\begin{table}[htbp]
\centering
\caption{Các bộ dữ liệu phát hiện biển, OCR và ký tự đơn lẻ đã tải thực tế.}

{\footnotesize Nguồn: Kaggle \autocite{tanhphp2026vnplatedet,bomaich2022vnplateyolov7,topkek2024vnplateocr,dtkngan2024motorbikeocr,kaggle2022vncharplate}.}
\label{tab:dataset-ocr}
\small
\begin{tabular}{p{3.4cm} p{3.0cm} p{2.0cm} p{3.0cm}}
\toprule
\textbf{Bộ dữ liệu} & \textbf{Số ảnh (thật)} & \textbf{Định dạng} & \textbf{Ghi chú} \\
\midrule
plate\_tanhphp~\autocite{tanhphp2026vnplatedet} & 2.255 (2.998 box) & YOLO bbox, nc=1 & Dùng làm nguồn detect chính bổ sung \\
plate\_bomaich\_yolov7~\autocite{bomaich2022vnplateyolov7} & 498 (780 box) & YOLO bbox, split sẵn & Ảnh cảnh đầy đủ, dùng bổ sung \\
plate\_ocr\_topkek~\autocite{topkek2024vnplateocr} & 12.190 (6.643 crop thật cộng 5.547 sinh tổng hợp) & Crop cộng CSV Name/Label/Type & Nguồn OCR chính; cần tách riêng phần sinh tổng hợp khi đánh giá để không thổi phồng accuracy \\
motorbike\_ocr\_100~\autocite{dtkngan2024motorbikeocr} & 116 & Ảnh cộng chuỗi biển & Tập test cho biển xe máy 2 hàng, chỉ có ảnh ban ngày \\
char\_nguyenquanglinh~\autocite{kaggle2022vncharplate} & 1.839 & Thư mục theo lớp, ảnh 28$\times$28 & Lệch lớp nặng giữa các ký tự; cần augment mạnh cho ký tự hiếm \\
\bottomrule
\end{tabular}
\end{table}

Các kho này không kèm file LICENSE rõ ràng nên đề tài chỉ dùng nội bộ, không phát hành
lại.

\subsection{Phân Tích Màu Biển Số (HSV)}
\label{sec:hsv}

Để phục vụ mô đun phân loại màu biển (Chương~\ref{ch:ocr}), đề tài chạy phân tích HSV trên
599 crop biển thật (từ bộ topkek\_plate\_ocr), tách vùng nền sáng khỏi vùng chữ tối.
Kết quả ở Bảng~\ref{tab:hsv-color}.

\begin{table}[htbp]
\centering
\caption{Phân bố màu nền trên 599 crop biển số thực tế.}

{\footnotesize Phân tích HSV trên crop từ bộ topkek\_plate\_ocr~\autocite{topkek2024vnplateocr}.}
\label{tab:hsv-color}
\begin{tabular}{lrrp{4.5cm}}
\toprule
\textbf{Màu nền} & \textbf{Số crop} & \textbf{Tỷ lệ} & \textbf{Ý nghĩa biển số VN} \\
\midrule
Trắng & 372 & 62,1\% & Xe cá nhân \\
Vàng & 115 & 19,2\% & Xe kinh doanh/dịch vụ \\
Xanh dương & 57 & 9,5\% & Cơ quan nhà nước \\
Đỏ & 37 & 6,2\% & Quân đội \\
\bottomrule
\end{tabular}
\end{table}

Bốn màu tách thành các cụm hue và saturation riêng biệt, cho thấy dùng HSV kết hợp CLAHE
mà không cần CNN riêng cho màu là khả thi. Dữ liệu lệch nặng về màu trắng (62,1\%), phản
ánh tỷ lệ thực của bãi xe dân sự.

\subsection{Bộ Dữ Liệu Phân Loại Xe}

Mô đun phân loại xe (Chương~\ref{ch:phanloai}) dùng A1 cho bước loại xe (ô tô, xe máy, xe
tải) và B5~\autocite{researchprojects2022vehiclebodystyle} cho bước kiểu dáng xe con.

% =====================================================================
\section{Trùng Lặp Gần giữa Tập Huấn Luyện và Tập Kiểm Tra}
\label{sec:dedup}
% =====================================================================

Trước khi đánh giá detector, đề tài so khớp average hash giữa tập huấn luyện và tập kiểm tra
của A1: 230 trong 572 ảnh kiểm tra (khoảng 40\%) có ảnh huấn luyện gần trùng với khoảng cách
Hamming không quá 5. Nguyên nhân là A1 trích từ luồng video liên tục nên các khung kế tiếp
chỉ khác vài pixel; chia ngẫu nhiên theo khung khiến tập kiểm tra rò rỉ từ tập huấn luyện
(data leakage) và thổi phồng chỉ số.

Đề tài loại các ảnh này trước khi đánh giá chính thức, tác động lên mAP trình bày ở
Bảng~\ref{tab:det-dedup}. Với dữ liệu trích từ video, nên chia tập theo thời gian thay vì
ngẫu nhiên theo khung.
